%! Simplifying Rational Expressions — Unit 4, Lesson 4.1 — Exit Ticket (Answer Key)
\documentclass[10pt]{article}
\usepackage{algebra2-article}
\usepackage{algebra2-key}   % pulls in -boxes; do NOT also load -boxes
\newcommand{\TallMath}[1]{$\displaystyle #1\rule[-1.4em]{0pt}{3.2em}$}

\begin{document}

\pageheader{Unit 4, Lesson 4.1}{Exit Ticket --- Answer Key}
\namedateperiod

\vspace{0.08in}

No notes. Show the factored line for question 1.

\vspace{0.06in}

\begin{enumerate}[leftmargin=1.8em, itemsep=9pt]
  \item \textbf{Simplify and restrict.} \; $\dfrac{x^2-2x-15}{x^2-9}$

        \vspace{4pt}
        Factored form: \ans{$\frac{(x-5)(x+3)}{(x-3)(x+3)}$}

        \vspace{4pt}
        Simplest form: \ans{$\frac{x-5}{x-3}$} \quad Restrictions: $x\neq$ \ans{$3,\ -3$}

        \vspace{4pt}
        One restriction is a \emph{hole} and one is a \emph{vertical asymptote}. \\[3pt]
        Hole at $x=$ \ans{$-3$} \quad Vertical asymptote at $x=$ \ans{3}

  \item \textbf{Explain.} A classmate says: ``Once the $(x+3)$ cancels, the expression is just
        $\dfrac{x-5}{x-3}$, so $x=-3$ is fine now.'' Is the classmate right? Explain what cancelling
        does and does not change.
        \ansline{No. Cancelling changes how the expression \emph{looks}, not which inputs it accepts.}
        \ansline{At $x=-3$ the original denominator $x^2-9$ is zero, so there is no value there to rewrite.}
        \ansline{The two forms agree at every \emph{other} input: the graph is $\frac{x-5}{x-3}$ with a hole at $\left(-3,\frac43\right)$.}

  \item \textbf{SOL-style multiple choice.} Which is the correct simplest form of
        $\dfrac{4-x^2}{x^2-x-2}$, with its restrictions?
        \begin{enumerate}[label=\Alph*., leftmargin=1.9em, itemsep=1pt]
          \item $\dfrac{x+2}{x+1}$, \; $x\neq2,\,-1$
          \item $-\dfrac{x+2}{x+1}$, \; $x\neq2,\,-1$
                \quad \textcolor{keyred}{\textbf{$\leftarrow$ correct}}
          \item $-\dfrac{x+2}{x+1}$, \; $x\neq-1$
          \item $-\dfrac{4}{x+2}$, \; $x\neq-2$
        \end{enumerate}
        \vspace{2pt}
        {\small $\dfrac{4-x^2}{x^2-x-2}=\dfrac{-(x-2)(x+2)}{(x-2)(x+1)}=-\dfrac{x+2}{x+1}$, $x\neq2,-1$.
        \textbf{A} drops the factor of $-1$ from the opposite binomials, \textbf{C} reads the
        restrictions off the \emph{simplified} denominator (losing the hole at $x=2$), and \textbf{D}
        cancels the $x^2$ \emph{terms}.}
\end{enumerate}

\begin{teachernote}
Graded for completion --- ``mistakes happen, blanks don't.'' Sort into four piles: \textbf{(1) got it};
\textbf{(2) restriction slip} --- listed only $x\neq3$, reading the restrictions off the simplified
denominator; \textbf{(3) sign slip} --- chose \textbf{A}, missing the opposite binomials; \textbf{(4)
term-cancelling} --- chose \textbf{D}, which means the Warm-Up's $\frac{3+6}{3}$ argument did not land.
Pile 2 is the one that matters most: it is the identical error that manufactures extraneous solutions
in Lesson~4.6, so a re-teach now is cheaper than one then. Item~2 earns full credit only if the student
says both halves: the value at $x=-3$ does not exist in the \emph{original}, and the two forms do agree
everywhere else.
\end{teachernote}

\end{document}
